\documentclass[../main.tex]{subfiles}
\begin{document}

\section*{Tổng kết kết quả và đóng góp học thuật}
\label{section:5.1}
\addcontentsline{toc}{section}{Tổng kết kết quả và đóng góp học thuật}

Luận văn đã hoàn thành toàn bộ mục tiêu đề ra, giải quyết thành công bài toán lượng hóa mức độ hòa trộn hành vi của mã độc trên hệ điều hành Windows và xác lập ranh giới rủi ro của các bộ phát hiện bất thường không giám sát (UAD). Thông qua phân tích lý thuyết hình học và thực nghiệm đối chứng trên 11 tập dữ liệu, nghiên cứu đã làm rõ những hạn chế hiệu năng biên của các phương pháp truyền thống và đề xuất hướng đi phi tham số để giải quyết bài toán mã độc ngụy trang.

Nghiên cứu mang lại hai đóng góp khoa học chính cho lý thuyết phát hiện bất thường và thực tiễn an ninh hệ thống:

\begin{itemize}
    \item \textbf{Đóng góp 1: Thiết lập thước đo độ khó nội tại (DensPct\_Mahal).} Bằng cách kết hợp hình học Gaussian (khoảng cách Mahalanobis) và đối chiếu với hàm phân phối tích lũy, thước đo này lượng hóa mức độ "chui sâu" của tiến trình độc hại vào trung tâm dữ liệu nền. Khảo sát thực nghiệm khẳng định thước đo này có khả năng giải thích sự sụp đổ hiệu năng (hiện tượng Early Collapse) của các thuật toán truyền thống chính xác và minh bạch hơn thước đo Marginal Overlap (DDO) thông thường.
    \item \textbf{Đóng góp 2: Phát hiện và chứng minh "Lỗ hổng cục bộ" (Local Pockets).} Nghiên cứu khám phá sự tồn tại của các vùng rỗng cục bộ nằm ngay trong lõi của phân phối lành tính, nơi mã độc LotL lẩn trốn sự giám sát toàn cục. Việc triển khai mô hình KDE\_Joint hoạt động như một \textit{bằng chứng tồn tại (existence proof)} xác nhận rằng hành vi ngụy trang tinh vi nhất (như chiến dịch APT29) hoàn toàn có thể bị bóc trần khi chuyển đổi phương pháp luận đo lường sang các công cụ phi tham số khai thác cấu trúc lân cận.
\end{itemize}

\section*{Hạn chế và Hướng phát triển tương lai}
\label{section:5.2}
\addcontentsline{toc}{section}{Hạn chế và Hướng phát triển tương lai}

Bên cạnh các kết quả đạt được, nghiên cứu cũng ghi nhận một số giới hạn thực tiễn cần khắc phục: (1) Cỡ mẫu thực nghiệm với ground truth đầy đủ còn tương đối nhỏ ($N=11$), làm giới hạn mức độ tổng quát của các kiểm định thống kê; (2) Công cụ KDE\_Joint sụp đổ hiệu năng trên không gian đặc trưng thưa thớt (sparse features), đòi hỏi sự can thiệp thủ công; (3) Môi trường thực nghiệm chưa xem xét hiện tượng trôi dạt khái niệm (Concept Drift) theo thời gian của hệ điều hành.

Từ những hạn chế trên, hướng phát triển trong tương lai được đề xuất như sau:

\begin{itemize}
    \item \textbf{Mở rộng quy mô đánh giá ($N \ge 30$):} Tích hợp các bộ dữ liệu an ninh quy mô công nghiệp (như DARPA OpTC, CICIDS) để làm phong phú bộ tham số kiểm định hoán vị, qua đó củng cố sức mạnh tổng quát hóa của khung đo lường.
    \item \textbf{Thiết kế Kiến trúc Hybrid Estimator (k-NN kết hợp KDE):} Để khắc phục giới hạn về chi phí thời gian $O(N^2 \cdot d)$ của KDE, một kiến trúc lai hai giai đoạn được đề xuất. Giai đoạn sàng lọc sử dụng cấu trúc cây kết hợp khoảng cách Gower ($O(N \log N)$) để tìm tập ứng viên láng giềng. Giai đoạn tiếp theo chỉ tính toán KDE trên tập láng giềng nhỏ ($O(k \cdot d)$). Thiết kế này hạ tổng chi phí vận hành xuống mức khả thi để áp dụng cho Trung tâm Điều hành An ninh (SOC) thời gian thực.
    \item \textbf{Phát triển Thước đo Local Pocket độc lập:} Nghiên cứu định nghĩa một \textit{Pocket Score} cấp độ dữ liệu, so sánh trực tiếp percentile toàn cục (Mahalanobis) và percentile cục bộ (k-NN) để lượng hóa mật độ cục bộ thành một đặc tính nội tại cấu trúc, từ đó đánh giá mức độ thưa thớt của mã độc mà không phụ thuộc vào kết quả của mô hình KDE.
    \item \textbf{Thích ứng với Concept Drift:} Tích hợp cơ chế cập nhật ma trận hiệp phương sai trực tuyến (\textit{Incremental Covariance Updating}) để tự động điều chỉnh khoảng cách Mahalanobis khi phân phối hành vi nền của hệ thống thay đổi, đảm bảo tính bền vững của thước đo theo thời gian.
\end{itemize}

\end{document}
